\chapter{ระบบหน่วยและการวัดทางฟิสิกส์}
\label{ch:ch01_units_measurement}
\index{ระบบหน่วยสากล SI}

\begin{conceptbox}[title=วัตถุประสงค์การเรียนรู้ประจำบท]
เมื่อสิ้นสุดการศึกษาบทเรียนนี้ นักศึกษาสามารถ
\begin{enumerate}
    \item อธิบายความหมายของปริมาณฐาน (base quantities) และปริมาณอนุพัทธ์ (derived quantities) พร้อมระบุปริมาณฐานทั้งเจ็ดของระบบหน่วยสากล (SI) ได้ถูกต้อง
    \item อธิบายระบบหน่วยสากล (International System of Units, SI) และความสำคัญของการมีระบบหน่วยมาตรฐานร่วมกันในวงการวิทยาศาสตร์และวิศวกรรม
    \item แปลงหน่วยระหว่างระบบต่าง ๆ และใช้คำอุปสรรค (prefixes) ในการเขียนและคำนวณปริมาณทางฟิสิกส์ได้อย่างถูกต้อง
    \item ใช้กฎเลขนัยสำคัญ (significant figures) ในการบันทึกผลการวัดและการคำนวณได้อย่างเหมาะสม
    \item แยกความแตกต่างระหว่างความแม่นยำ (precision) กับความถูกต้อง (accuracy) ของการวัด
\end{enumerate}
\end{conceptbox}

\section{ทฤษฎีและหลักการพื้นฐาน}

\begin{bioappbox}
\textbf{การวัดและการแปลงหน่วยในงานจุลชีววิทยาและวิทยาศาสตร์สิ่งแวดล้อม}

ในงานจุลชีววิทยาและวิทยาศาสตร์สิ่งแวดล้อม การวัดขนาดของวัตถุที่มีขนาดเล็กมากระดับไมโครเมตร ($\mu\text{m}$) และนาโนเมตร ($\text{nm}$) มีความสำคัญยิ่ง เช่น การวัดขนาดของแบคทีเรีย \textit{Escherichia coli} ซึ่งมีความยาวเฉลี่ยประมาณ $2.0\,\mu\text{m}$ ($2.0 \times 10^{-6}\,\text{m}$) และเส้นผ่านศูนย์กลางของอนุภาคฝุ่นละออง PM2.5 ที่มีขนาดไม่เกิน $2.5\,\mu\text{m}$ นักศึกษาจำเป็นต้องมีความเชี่ยวชาญในการแปลงหน่วยตามระบบเอสไอและการระบุเลขนัยสำคัญ เพื่อให้ผลการวิเคราะห์มีความแม่นยำและน่าเชื่อถือตามเกณฑ์มาตรฐานสากล
\end{bioappbox}

\begin{labbox}
1. ตรวจสอบตำแหน่งศูนย์ (Zero Error) ของเครื่องมือวัดทุกชนิดก่อนเริ่มบันทึกข้อมูล
2. ทำการวัดซ้ำอย่างน้อย 3 ถึง 5 ครั้งในแต่ละจุดทดลองเพื่อคำนวณหาค่าเฉลี่ยและค่าเบี่ยงเบนมาตรฐาน
3. บันทึกผลการวัดพร้อมระบุหน่วยเอสไอ (SI Units) และค่าความคลาดเคลื่อนของการวัดเสมอ
\end{labbox}

\section{ตัวอย่างโจทย์และการวิเคราะห์เชิงฟิสิกส์}

\begin{examplebox}
ตัวอย่างที่ 1.1

รถยนต์คันหนึ่งเคลื่อนที่ด้วยอัตราเร็ว 90 กิโลเมตรต่อชั่วโมง (km/hr) จงแปลงอัตราเร็วนี้ให้อยู่ในหน่วยเมตรต่อวินาที (m/s)

วิธีทำ

ใช้ factor-label method โดยแปลงหน่วยความยาวจาก km เป็น m และหน่วยเวลาจาก hr เป็น s
\end{examplebox}

\section{แบบฝึกหัดทบทวนและคำถามท้ายบท}

\begin{enumerate}
    \item จงอธิบายความแตกต่างระหว่างปริมาณฐานและปริมาณอนุพัทธ์ พร้อมยกตัวอย่างอย่างละสองปริมาณ
    \item จงระบุปริมาณฐานทั้งเจ็ดของระบบ SI พร้อมหน่วยและสัญลักษณ์
    \item คำอุปสรรค kilo, mega และ giga แทนตัวพหุคูณของสิบเท่าใดตามลำดับ
    \item จงแปลง 5.0 กิโลเมตร ให้อยู่ในหน่วยเซนติเมตร
    \item จงแปลง 250 มิลลิลิตร ให้อยู่ในหน่วยลิตร (กำหนดให้ 1 ลิตร = 1000 มิลลิลิตร)
    \item รถไฟความเร็วสูงเคลื่อนที่ด้วยอัตราเร็ว 250 km/hr จงแปลงเป็นหน่วย m/s
\end{enumerate}

% สิ้นสุดบทเรียน
